\chapter{Justificativa e Histórico}

O curso de Engenharia de Computação da Universidade Federal de Uberlândia, desde sua criação como um curso independente\footnote{A FEELT já trazia a marca da engenharia da computação desde 1999, com o certificado disponível para os discentes de Engenharia Elétrica.} em 2012~\cite{UFU:CONSUN:20:2012}, tem se destacado pela formação de profissionais de alta qualificação. Sua trajetória é marcada pela busca constante por excelência e atualização, refletida no processo de modernização curricular iniciado em 2018, que resultou na implementação do Projeto Pedagógico de Curso (PPC) da versão 2019-1~\cite{UFU:FEELT:PPC:EC:2019}. Este PPC foi um passo fundamental para adequar o curso às demandas contemporâneas da época, em especial o alinhamento às Diretrizes Curriculares Nacionais (DCNs) da área de Computação~\cite{MEC:CNE:CES:5:2016}.

Contudo, a velocidade das transformações tecnológicas e a crescente complexidade dos desafios globais exigem que a formação em engenharia avance continuamente. É com este propósito que se apresenta esta nova versão do PPC para 2025. A atualização fundamenta-se em três pilares estratégicos: a adoção de referenciais curriculares de ponta, como o \textit{Computing Curricula 2020} (CC2020) da ACM~\cite{ACM:CC2020}; a integração de competências transversais de sustentabilidade, baseadas nos princípios do Engineering Education for Sustainable Development (EESD)~\cite{UNESCO:Eng4SusDev:2021}; e a modernização das práticas pedagógicas através do ensino híbrido~\cite{Decreto:12456:2025}. Juntas, essas diretrizes visam fortalecer a formação de nossos estudantes, com o objetivo de entregar à sociedade egressos com um perfil mais internacionalizado e com maior capacidade de gerar impacto positivo, seja na indústria, na pesquisa ou no empreendedorismo.

\section{Justificativa}

Nos últimos anos, a expansão acelerada das tecnologias digitais, a convergência entre hardware e software, o crescimento de áreas como inteligência artificial, sistemas embarcados, computação em nuvem, internet das coisas e segurança cibernética impuseram novos desafios à formação em Engenharia de Computação. Ao mesmo tempo, a sociedade exige de seus profissionais um perfil mais ético, comunicativo, colaborativo e sensível às realidades socioculturais em que estão inseridos.

Adicionalmente, este PPC responde ao chamado das políticas públicas e institucionais que estabelecem novos compromissos para os cursos de graduação. Destaca-se, nesse sentido, a obrigatoriedade da curricularização da extensão (conforme a meta 12.7 do Plano Nacional de Educação), a valorização da internacionalização, a possibilidade de inserção de componentes híbridos ou em EaD (até 30\% da carga horária), e a adoção de modelos formativos centrados em competências -- como preconizado pelo CC2020 da ACM/IEEE que indica a adoção de \textit{Competency-Based Learning} (CBL). A Universidade Federal de Uberlândia, por meio do seu Plano Institucional de Desenvolvimento e Expansão (PIDE)~\cite{UFU:CONSUN:31:2022} e das resoluções do CONGRAD, orienta-se por essas diretrizes ao estimular a inovação pedagógica, a interdisciplinaridade, a equidade e a responsabilidade social.

A reestruturação do curso em 2025 também parte de um olhar atento para o contexto regional do Triângulo Mineiro e do interior do Brasil. A vocação tecnológica da região, os ecossistemas de inovação em expansão, os arranjos produtivos locais e a atuação crescente de startups, empresas de base tecnológica, laboratórios públicos e instituições de ensino e pesquisa demandam uma formação versátil, empreendedora e conectada com os desafios locais e globais.

Assim, a atualização deste PPC visa não apenas cumprir exigências normativas, mas reafirmar o compromisso do curso com a excelência acadêmica, a pertinência social e a formação integral do engenheiro de computação, preparado para intervir criticamente nos mais diversos campos da tecnologia, da ciência e da sociedade.

\subsection{Justificativa para Alteração da Denominação do Curso}

A alteração da denominação do curso de graduação em ``Engenharia de Computação'' para ``\ppccurso'' fundamenta-se em um conjunto de razões acadêmicas, institucionais e sociais que refletem a evolução do curso, a aderência às políticas nacionais de educação superior e a pertinência de sua atuação frente às transformações tecnológicas contemporâneas.

% Do ponto de vista acadêmico e científico, a Inteligência Artificial deixou de ser um campo periférico da computação para consolidar-se como um eixo de referência, transversal e estratégico. O \textit{Computing Curricula 2020}~\cite{ACM:CC2020} reconhece a IA como parte essencial da formação em Computação, destacando-a como elemento de integração entre hardware, software, ciência de dados e sistemas inteligentes. Ao incorporá-la na própria denominação, o curso reafirma sua atualização frente às referências internacionais e sua vocação para formar engenheiros capazes de atuar nas fronteiras do conhecimento tecnológico, sem perder a tradição de excelência na integração entre eletrônica, sistemas digitais, redes e software.

Do ponto de vista acadêmico e científico, a Inteligência Artificial deixou de ser um campo periférico da computação para consolidar-se como uma área de especialização estratégica e transversal. O \textit{Computing Curricula 2020}~\cite{ACM:CC2020} reconhece a IA como parte essencial da formação em Computação, destacando-a como elemento de integração entre hardware, software, ciência de dados e sistemas inteligentes. Ao incorporá-la na própria denominação, o curso reafirma sua atualização frente às referências internacionais e sua vocação para formar engenheiros capazes de atuar nas fronteiras do conhecimento tecnológico, sem perder a tradição de excelência na integração entre eletrônica, sistemas digitais, redes e software.

% Além da mudança de nome, a reformulação curricular trouxe um aumento expressivo de componentes obrigatórios em Inteligência Artificial, com carga horária quatro vezes superior à anteriormente dedicada à área, além da possibilidade de ampliação por meio de disciplinas optativas. Esse avanço assegura a todos os(as) estudantes uma formação sólida e progressiva em fundamentos, técnicas e aplicações de IA. Mais do que isso, a IA não se restringe às disciplinas específicas: o PPC prevê seu uso ético, crítico e consciente no processo de ensino-aprendizagem, articulando formação técnica, reflexão ética e práticas pedagógicas inovadoras (ver \autoref{sec:IA-UFU}). Assim, a Inteligência Artificial consolida-se tanto como componente curricular diferenciador e aplicado quanto como recurso pedagógico transversal, desenvolvendo nos discentes competências de letramento digital, autonomia e responsabilidade no uso de tecnologias emergentes.

Além da mudança de nome, a reformulação curricular trouxe um aumento expressivo de componentes obrigatórios em Inteligência Artificial, com carga horária quatro vezes superior à anteriormente dedicada à área, além da possibilidade de ampliação por meio de disciplinas optativas. Essa expansão ocorre sem prejuízo da formação tradicional e consolidada em Engenharia de Computação, que mantém sua robustez em áreas como eletrônica, sistemas embarcados, arquitetura de computadores, redes de comunicação e sistemas de controle — fundamentos essenciais para a aplicação eficaz de técnicas de IA em sistemas computacionais reais. O termo ``Aplicada'' na denominação reforça a vocação do curso para formar engenheiros capazes de aplicar técnicas de IA na resolução de problemas reais de engenharia, tais como cibersegurança, sistemas embarcados inteligentes, otimização de redes, processamento de sinais e sistemas de controle adaptativos. Esse avanço assegura a todos os(as) estudantes uma formação sólida e progressiva em fundamentos, técnicas e aplicações de IA. Mais do que isso, a IA não se restringe às disciplinas específicas: o PPC prevê seu uso ético, crítico e consciente no processo de ensino-aprendizagem, articulando formação técnica, reflexão ética e práticas pedagógicas inovadoras (ver \autoref{sec:IA-UFU}). Assim, a Inteligência Artificial consolida-se tanto como componente curricular diferenciador e aplicado quanto como recurso pedagógico transversal, desenvolvendo nos discentes competências de letramento digital, autonomia e responsabilidade no uso de tecnologias emergentes.

Esse fortalecimento interno dialoga diretamente com movimentos mais amplos em nível nacional. Sob a perspectiva institucional e de políticas públicas, a alteração da denominação alinha-se ao Plano Brasileiro de Inteligência Artificial (PBIA)\footnote{\href{https://instituto.ia.lncc.br/pt/noticias/versao-final-plano-brasileiro-de-inteligencia-artificial-pbia}{Plano Brasileiro de Inteligência Artificial (PBIA), junho/2025}} e ao programa Universidades Inovadoras e Sustentáveis\footnote{\href{https://www.gov.br/fazenda/pt-br/assuntos/noticias/2025/outubro/programa-universidades-inovadoras-e-sustentaveis-e-lancado-em-parceria-com-o-mec-no-festival-curicaca}{Programa Universidades Inovadoras e Sustentáveis, outubro/2025}}, ambos do Governo Federal, que definem a expansão da formação em IA como prioridade estratégica. O PBIA estabelece metas até 2028 e 2030, como a criação de 5.000 novas vagas de graduação em IA e o aumento em 50\% do número de formados em áreas STEM com ênfase em inteligência artificial. Nesse cenário, a UFU posiciona-se de forma protagonista ao atualizar a identidade do curso, consolidando sua relevância no contexto do ensino superior público brasileiro.

No plano regional e social, o Triângulo Mineiro e o interior do Brasil constituem ecossistemas de inovação em expansão, com startups, empresas de base tecnológica, laboratórios públicos e centros de pesquisa aplicados. Em particular, o lançamento do programa Uberlând.I.A. pela Prefeitura de Uberlândia\footnote{\href{https://www.uberlandia.mg.gov.br/2025/09/04/prefeitura-promove-lancamento-do-programa-uberland-i-a/}{Programa Uberlând.I.A., setembro/2025}} — que reúne instituições de ensino, empresas de tecnologia e poder público para fortalecer a cidade como protagonista da era da IA — evidencia a oportunidade de articulação direta entre o curso, os agentes regionais de inovação e o mercado local. A presença da Inteligência Artificial no nome do curso amplia sua atratividade para novos estudantes, reforça sua liderança na formação de mão de obra qualificada e fortalece parcerias com setores estratégicos como saúde, agricultura, logística, energia e transformação digital. Essa atualização responde, portanto, não apenas a demandas acadêmicas, mas também a demandas sociais e econômicas, reafirmando o compromisso da UFU com o desenvolvimento regional e nacional.

% Por fim, a alteração da denominação contribui para a visibilidade e a internacionalização do curso. Em um cenário em que programas de graduação em \textit{Computer Engineering with Artificial Intelligence} já são realidade em universidades estrangeiras, a mudança aproxima a UFU dos referenciais globais, facilita a cooperação acadêmica internacional e aumenta a competitividade de seus egressos no mercado de trabalho. Mais do que um ajuste nominal, a nova denominação traduz a identidade contemporânea do curso, projetando-o para o futuro sem perder a solidez de sua trajetória.

Por fim, a alteração da denominação contribui para a visibilidade e a internacionalização do curso. Em um cenário em que programas de graduação que integram Engenharia de Computação e Inteligência Artificial já são realidade em universidades estrangeiras de referência, a mudança aproxima a UFU dos referenciais globais sem perder as especificidades e fortalezas regionais do curso, facilita a cooperação acadêmica internacional e aumenta a competitividade de seus egressos no mercado de trabalho. Mais do que um ajuste nominal, a nova denominação traduz a identidade contemporânea do curso, projetando-o para o futuro sem perder a solidez de sua trajetória.

\subsection{Justificativa para a Ampliação do Número de Vagas}

A proposta de alteração do Projeto Pedagógico do Curso (PPC) de \ppccurso (anteriormente denominado Engenharia de Computação) contempla a ampliação do número de vagas semestrais de 15 (quinze) para 30 (trinta). Tal medida se fundamenta em uma análise robusta que demonstra uma demanda consolidada e crescente, uma eficiência ímpar na gestão de vagas, viabilidade operacional e total alinhamento com a missão estratégica da Universidade Federal de Uberlândia (UFU).

\subsubsection*{Demanda Externa Consolidada e Crescente}

O curso de Engenharia de Computação (atualmente denominado \ppccurso) demonstra consistentemente uma demanda muito superior à sua oferta. Mesmo sendo responsável por apenas 16,7\% das vagas de entrada da Faculdade de Engenharia Elétrica (FEELT), o curso atrai, em média, \textbf{30\% de todos os candidatos da faculdade via SISU e 50\% via Vestibular}, conforme dados dos últimos quatro anos. Isso resulta em uma das maiores relações candidato/vaga da universidade e extensas listas de espera.

Além da alta procura no ingresso, o curso atende a uma forte demanda do mercado de trabalho, inserido em um dos principais polos tecnológicos do Brasil. A alta taxa de empregabilidade dos egressos e a crescente procura por estagiários validam a necessidade de formar mais profissionais para suprir as demandas do arranjo produtivo local e nacional.

\subsubsection*{Sustentabilidade e Eficiência na Gestão de Vagas}

Em um contexto de otimização de recursos públicos, a Engenharia de Computação (atualmente denominado \ppccurso) se destaca como um modelo de sustentabilidade. Conforme levantamento institucional de 2022-2025, enquanto a FEELT enfrentou perda de vagas discentes, como muitas outras Unidades Acadêmicas desta Universidade, este curso foi responsável por apenas 0,69\% dessas perdas.

A eficiência do curso também se manifesta na sua capacidade de reter o valor social de suas vagas. Com uma taxa de reaproveitamento de vagas ociosas de praticamente 100\%, o curso atraiu em média\footnote{Nos últimos quatro anos (2022-2025)} 3,26 estudantes para cada vaga que surgiu por desistência, jubilamento ou não preenchimento em processos seletivos (vagas reservadas específicas). Isso comprova uma forte demanda reprimida e garante que o investimento público cumpra seu papel social, em nítido contraste com outros cursos da unidade.

\subsubsection*{Demanda Interna e Viabilidade Operacional}

A ampliação é plenamente viável. A análise de ocupação de disciplinas demonstra que a \ppccurso é o único curso da FEELT que opera com superávit de matrículas, superando sua própria expectativa em mais de 30\%\footnote{No semestre letivo de 2025/1}, o que indica uma intensa procura por suas disciplinas por alunos de outros cursos.

O aumento proposto é absorvível pela estrutura existente. A taxa média de vagas ofertadas e não utilizadas nas disciplinas da FEELT é de \textbf{36\%}\footnote{Média dos semestres letivos de 2024/1 a 2025/1}, um espaço que pode acomodar o crescimento. O dobro de ingressantes na Engenharia de Computação representa um impacto gerenciável de apenas \textbf{13,2\%} na expectativa total de matrículas da faculdade. As turmas específicas passariam a ter uma média de 30 alunos, e os discentes adicionais nas disciplinas básicas ajudariam a mitigar a disponibilidade gerada por outros cursos.

\subsubsection*{Alinhamento Estratégico com a Missão da UFU}

A ampliação de vagas está em plena consonância com os objetivos do Plano Institucional de Desenvolvimento e Expansão (PIDE), contribuindo para:

\begin{itemize}
    \item Ampliar o acesso a um curso de alta relevância social e tecnológica.
    \item Fortalecer o papel da Universidade no desenvolvimento do ecossistema de inovação regional.
    \item Otimizar o uso de recursos, por meio de uma decisão de gestão estratégica que realoca o investimento público de vagas cronicamente ociosas para um curso de alta performance, eficiência e comprovado retorno social e acadêmico.
\end{itemize}

Diante do exposto, a ampliação do número de vagas é uma ação estratégica, justificada pela demanda externa e interna, validada pela eficiência operacional e alinhada à missão da Universidade de servir à sociedade com ensino público de excelência.

\section{Histórico do Curso}

A trajetória do curso de Engenharia de Computação na Universidade Federal de Uberlândia (UFU) iniciou-se em 1999, quando foi criado como um Certificado vinculado ao curso de Engenharia Elétrica, por meio da Resolução nº 4/1999 do Conselho Universitário. As atividades acadêmicas desta primeira fase começaram no ano 2000.

Um novo marco foi estabelecido em 2012, quando o curso foi instituído como uma graduação independente~\cite{UFU:CONSUN:20:2012}, culminando na adoção de seu primeiro Projeto Pedagógico de Curso e no ingresso da primeira turma em 2013. Desde sua origem, o curso consolidou uma proposta pedagógica inovadora, unindo uma sólida formação em fundamentos matemáticos e computacionais com o aprofundamento técnico nas áreas de eletrônica, sistemas digitais e engenharia de software.

Nas primeiras décadas de funcionamento, o curso destacou-se pelo pioneirismo na integração entre hardware e software em uma perspectiva sistêmica, pela intensa atividade de pesquisa e extensão conduzida por docentes altamente qualificados e pelo envolvimento dos discentes em projetos de iniciação científica, inovação e empreendedorismo tecnológico. A infraestrutura laboratorial e o corpo docente vinculado à Faculdade de Engenharia Elétrica (FEELT) tornaram-se pilares do desenvolvimento acadêmico do curso, que gradualmente passou a ocupar posição de destaque entre os cursos de Engenharia da UFU.

A primeira grande reformulação curricular ocorreu em 2018, com foco na atualização de conteúdos, na adequação da carga horária às novas diretrizes da UFU e na consolidação das disciplinas de programação, arquitetura de computadores e sistemas embarcados. Uma nova atualização curricular foi implementada, em consonância com a Resolução CNE/CES nº 5/2016, que redefiniu as Diretrizes Curriculares Nacionais para os cursos da área de Computação. Esse PPC incorporou, de forma mais orgânica, as atividades de extensão, reforçou a formação empreendedora, introduziu disciplinas optativas com foco em tecnologias emergentes e aperfeiçoou os componentes de estágio e trabalho de conclusão de curso (TCC).

A partir de 2020, o curso enfrentou os impactos da pandemia de COVID-19, que exigiram adaptações rápidas em sua organização didático-pedagógica. A experiência com o ensino remoto emergencial fortaleceu a capacidade do curso de explorar novas metodologias de ensino-aprendizagem, potencializando o uso de tecnologias digitais, práticas híbridas e estratégias de avaliação diversificadas. Esse período também impulsionou reflexões sobre a necessidade de maior flexibilidade curricular, acolhimento aos discentes e revisão de objetivos formativos.

Em 2026, a \ppccurso chega a uma nova etapa de sua trajetória com a presente atualização do Projeto Pedagógico. Essa reformulação reafirma os princípios fundadores do curso — excelência técnica, formação crítica, compromisso social — e os amplia à luz das demandas contemporâneas: formação orientada por competências, interdisciplinaridade, internacionalização, inclusão, inovação pedagógica e articulação com os ecossistemas tecnológicos da região e do país.

\section{Unidade Acadêmica}

O curso de \ppccurso está vinculado à Faculdade de Engenharia Elétrica (FEELT) da Universidade Federal de Uberlândia (UFU), unidade acadêmica responsável por sua gestão pedagógica, administrativa e acadêmica. A FEELT atualmente abriga seis cursos de graduação presenciais — cinco localizados no campus Santa Mônica, em Uberlândia/MG (Engenharia Elétrica, Engenharia Biomédica, Engenharia de Computação, Engenharia Eletrônica e de Telecomunicações, Engenharia de Controle e Automação) e um no campus Patos de Minas/MG (Engenharia Eletrônica e de Telecomunicações) — além de dois programas de pós-graduação stricto sensu: o Programa de Pós-Graduação em Engenharia Elétrica e o Programa de Pós-Graduação em Engenharia Biomédica.

Com a aprovação do novo Regimento Interno da FEELT~\cite{UFU:FEELT:Regimento:2023}, a estrutura organizacional da unidade foi reorganizada em Departamentos Acadêmicos, que substituíram os antigos núcleos temáticos. Cada departamento é responsável por atividades de ensino, pesquisa e extensão em áreas específicas do conhecimento, organizando a atuação docente e a oferta de disciplinas para os diversos cursos de graduação e pós-graduação.

Os departamentos atualmente constituídos na FEELT são:

\begin{itemize}
    \item Departamento de Sistemas de Potência e Energia (DSPE)
    \item Departamento de Engenharia Biomédica (DEB)
    \item Departamento de Computação e Inovação em Engenharia (DCIE)
    \item Departamento de Sistemas de Controle e Automação (DSCA)
    \item Departamento de Eletrônica e Telecomunicações (DETEL)
    \item Departamento de Máquinas Elétricas e Materiais (DMEM)
\end{itemize}

O curso de \ppccurso tem sua base acadêmica fortemente integrada ao Departamento de Computação e Inovação em Engenharia (DCIE), embora conte com a contribuição de docentes de outros departamentos da FEELT, bem como de unidades acadêmicas externas, como a Faculdade de Matemática (FAMAT), o Instituto de Física (INFIS) e a Faculdade de Computação (FACOM), na oferta de disciplinas obrigatórias e optativas.

A estrutura da FEELT também favorece a participação dos(as) discentes em programas de formação complementar, como mobilidade acadêmica nacional e internacional, Empresa Júnior (CONSELT), Programa de Educação Tutorial (PET), monitoria, iniciação científica (PIBIC, PIVIC/UFU), projetos de extensão e programas de cooperação acadêmica (ex: CAPES/Brafitec). No âmbito da pós-graduação, os(as) docentes da FEELT atuam em linhas de pesquisa avançadas nas áreas de inteligência artificial, sistemas embarcados, processamento de sinais, realidade virtual, bioengenharia, fontes alternativas de energia, entre outras.

Essa estrutura acadêmica consolidada e integrada garante ao curso de \ppccurso uma base sólida para a formação de profissionais críticos, inovadores e aptos a atuar nos mais diversos contextos tecnológicos e sociais.

\section{Relação entre Sociedade e o Curso}

Desde sua criação, o curso de Engenharia de Computação da Universidade Federal de Uberlândia tem se orientado pela missão de formar profissionais capazes de responder às necessidades da sociedade contemporânea, contribuindo para o avanço científico, tecnológico, econômico e social da região e do país. Sua atuação se articula com o contexto local e regional, integrando-se ao ecossistema de inovação do Triângulo Mineiro e ampliando sua inserção em redes nacionais e internacionais de cooperação.

A cidade de Uberlândia/MG, onde o curso está localizado, destaca-se como um dos principais polos tecnológicos do interior do Brasil, combinando infraestrutura logística estratégica, mão-de-obra qualificada e ambiente favorável ao desenvolvimento de negócios de base tecnológica. A presença de backbones digitais, redes de fibra óptica, cobertura 5G, centros de dados, hubs de inovação e conexão com os grandes centros urbanos — como São Paulo, Brasília, Goiânia e Belo Horizonte — fortalece o posicionamento da cidade como espaço propício à pesquisa aplicada, à inovação e ao empreendedorismo.

Nos últimos anos, diversas instituições e iniciativas locais vêm consolidando esse ecossistema, como a Secretaria Municipal de Desenvolvimento Econômico, Inovação e Turismo, a i9 Uberlândia, a Comunidade Colmeia, a Minas Startup, o Parque do Sabiá Tech, e o SEBRAE/MG. Essas organizações têm fomentado a articulação entre universidades, startups, empresas consolidadas e o poder público, promovendo ações de apoio à inovação, aceleração de negócios e incentivo à criação de soluções tecnológicas com impacto social e econômico.

Nesse contexto, o curso de Engenharia de Computação vem desempenhando papel estratégico, tanto pela formação de profissionais qualificados quanto pela sua produção científica e tecnológica. A atuação dos(as) docentes em projetos de pesquisa, desenvolvimento e inovação (PD\&I) tem resultado em parcerias com grandes empresas e instituições — como ANEEL, CEMIG, Algar Telecom, ARCOM — e em projetos financiados por agências de fomento nacionais (CAPES, CNPq, FINEP, FAPEMIG). A colaboração com empresas e centros tecnológicos estimula não apenas a pesquisa aplicada, mas também a inserção precoce dos(as) estudantes em experiências reais de engenharia.

No plano internacional, o curso mantém convênios ativos, especialmente por meio de programas como o Brafitec (CAPES), com ênfase na cooperação com instituições francesas, e busca ampliar sua presença global por meio de parcerias bilaterais, eventos acadêmicos e a recepção de estudantes estrangeiros. Ainda que modestamente, o número de discentes estrangeiros e as experiências de intercâmbio internacional vêm crescendo.

A formação oferecida no curso também tem impactado positivamente o programa de pós-graduação em Engenharia Elétrica da FEELT. Dados recentes indicam que uma parcela significativa das dissertações e teses defendidas na unidade concentra-se em temas associados à Engenharia de Computação, evidenciando a sinergia entre graduação e pós-graduação e a consolidação de uma comunidade acadêmica ativa e produtiva.

A atualização deste PPC reafirma esse compromisso com a sociedade, respondendo à crescente demanda por educação pública de qualidade, voltada para áreas de alta relevância científica e tecnológica, e promovendo uma formação alinhada às exigências contemporâneas da cidadania, da sustentabilidade e da inovação.
